\section{Collection and document management: organisation, classification and
lifecycle}
\label{sec:organization}

Collection membership defines the scope available to a query. HMRS uses
the subject of the content for placement, preserves the source assignment
across split parts, and exposes descriptions before an agent opens a
document body. The collection and document operations in
Table~\ref{tab:levels} maintain these choices over the life of the corpus.

\subsection{Classification, metadata, tags and duplicate review}
\label{sec:org-placement}
\label{sec:org-registries}
\label{sec:org-tags}
\label{sec:org-dedup}

The catalogue returns collections and sub-libraries with descriptions.
Each shelf shows document titles and descriptions within its selected scope.
The ingest agent uses that scope and the parsed content to select a library,
then writes descriptions and tags grounded in the body. Duplicate review
presents candidate pairs for inspection because adjacent parts may share
metadata while retaining distinct source sections.

HMRS supports library and sub-library creation, metadata edits, moves,
renames, deletion and forced re-indexing. A move updates the document path
and registry entries, while re-indexing updates the searchable representation.
Integrity checks compare document existence, path and index metadata across
the files on disk, the \texttt{.tree-fs.json} registry and the library's
\texttt{.knowledge-base.yml}. These checks give operators observable
outcomes for the state changes made by each operation.

Sub-libraries keep their descriptions and indices at their own tree nodes.
When the caller supplies a parent scope, the vector service resolves the
parent and its descendants into the collections searched. The lexical stage
of the two-stage service likewise resolves descendant identifiers. The
tested flat union calls explicit library scopes, which fixes the collections
compared in the evaluation.

\subsection{Lifecycle operations, graph and operational knowledge}
\label{sec:org-lifecycle}
\label{sec:org-graph}
\label{sec:org-experience}

One historical rename left old-path chunks until a forced re-index. The
current rename paths attempt to synchronise index entries and report
failures separately from the metadata update. The post-operation query in
E2 tests the result in one fixture, not the general invariant.

The document graph links libraries, documents and tags for navigation and
candidate expansion. Operational lessons are also stored as retrievable
documents. Experience tools track reviews, use and freshness and use
credibility tiers to order results. Section~\ref{sec:eval} reports
operational-experience cases alongside the management experiments.

\subsection{Organisation contract}
\label{sec:org-integrity}
\label{sec:org-summary}

The organisation contract links each document to its current owning
collection, descriptive metadata and indexed path. The checks in
Section~\ref{sec:e2} test these relationships on selected lifecycle
operations by reconciling the three sources, reviewing document descriptions
and rerunning retrieval probes. The resulting traces locate inconsistencies
at the file, registry or index stage.
